\subsection{Proper canonical point fields}
\label{sec:proper-canonical-point-fields}
An exchanged canonical line strengthens the point-field degree
constraints at proper torsion levels. The argument uses an exact lower
break only when a coordinate has valuation prime to $p$ in its
actual point field.

Pass to the completed maximal unramified base and take the common
tame good field $L/K$ of degree $p^2-1$, with algebraically closed
residue field and $v_L(L^\times)=\Z$. Write $e=e_K$.

\begin{proposition}\label{prop:proper-canonical-point-fields}
Let $p\ge5$, and suppose two good models have the same central
parameter $e<z<pe$. Put
\[
 A=(p+1)e-z,\qquad B=z/p,\qquad a=v_p(A)\ge1.
\]
Let $f:G[p^n]\simeq H[p^n]$, $n\ge3$, be an actual torsion
isomorphism carrying the canonical line $C_G$ to a noncanonical
line of $H$. For an exact $p^j$-point $Q$ with
$p^{j-1}Q\in C_G\setminus\{0\}$, its point-field degree is a
power $[L(Q):L]=p^d$. If $2\le j<n$, then
\[
 d\ge\max\{0,2j-1-a\}.
\]
At $j=n$ one always has $d\ge\max\{0,2n-2-a\}$.
If $d=2n-2-a\ge0$, then $L(Q)/L(pQ)$ has degree at most $p$,
and its translation subgroup has trivial intersection with $C_G$.
In particular a nonzero such subgroup is a noncanonical line.
\end{proposition}

\begin{proof}
The canonical points of each group are $L$-rational by
Lemma~\ref{fc-lem:rational}. The image $f(C_G)$ is also pointwise
$L$-rational and differs from $C_H$. These two lines span $H[p]$,
so $H[p]$, and then $G[p]$, are pointwise $L$-rational. In
particular $B$ is an integer. No higher cyclic subgroup has been
inferred from this argument.

Choose any compatible tower $Q_1,\ldots,Q_n$ with
$pQ_i=Q_{i-1}$ and $Q_1\in C_G\setminus\{0\}$, and put
$R_i=f(Q_i)$. Since $A<p^2B$, the Newton division formula gives
\begin{equation}\label{eq:proper-point-values}
 v_L(T(Q_i))=\frac{A}{p^{2i-2}},\qquad
 v_L(S(R_i))=\frac{B}{p^{2i-2}}.
\end{equation}
For the first formula the initial target is canonical and lies
strictly below the middle division breakpoint; later target
values decrease. The second begins at a noncanonical point and
always remains on the degree-$p^2$ branch. Thus these values hold
for every choice of lifts, by Lemma~\ref{rc-lem:newton}.

Let $F_i=L(Q_i)$. The actual isomorphism identifies point
stabilizers, so $F_i=L(R_i)$. Each extension $F_i/F_{i-1}$ is
Galois: every conjugate of $Q_i$ over $F_{i-1}$ is $Q_i+P$ for
some $P\in G[p](L)$, and its coordinate belongs to $F_i$ by
formal addition. The map
\[
 \sigma\longmapsto\sigma(Q_i)-Q_i
\]
identifies its Galois group with an $\F_p$-subspace
$V_i\subset G[p](L)$. It is a homomorphism because all the
translating points are fixed over $L$. Consequently
\[
 [F_i:L]=p^{d_i},\quad d_1=0,\quad
 d_i-d_{i-1}=\dim_{\F_p}V_i\in\{0,1,2\},\quad d_i\le2i-2.
\]
The extensions are totally ramified because the residue field
is algebraically closed. No assertion that $F_i/L$ is Galois
has been used.

We first prove that
\begin{equation}\label{eq:prime-coordinate-no-canonical-translation}
 p\nmid w_i(T(Q_i))\quad\Longrightarrow\quad
 V_i\cap C_G=0\qquad(i\ge2),
\end{equation}
where $w_i=p^{d_i}v_L$ is the integer valuation of $F_i$.
Suppose instead that an automorphism $\sigma$ translates $Q_i$
by a nonzero point $P_0\in C_G$. It has order $p$. Apply
Lemma~\ref{cr-lem:valuation} to the cyclic extension of $F_i$
over its fixed field under $\sigma$.

For a good formal law, $F(T,U)-T$ is $U$ times a power series
with constant term one. Evaluated at points of positive value,
the remaining factor is a unit. Thus
\[
 w_i\bigl(\sigma(T(Q_i))-T(Q_i)\bigr)=p^{d_i}A.
\]
The prime-valuation clause of Lemma~\ref{cr-lem:valuation}
determines the lower break in this particular cyclic extension:
\[
 b_\sigma=p^{d_i}A\bigl(1-p^{-2i+2}\bigr).
\]
For the transported point the same automorphism gives
$\sigma(R_i)=R_i+f(P_0)$, and $f(P_0)$ is noncanonical, of
value $B$. The unconditional inequality of the same lemma and
\eqref{eq:proper-point-values} yield
\[
 b_\sigma\le p^{d_i}B\bigl(1-p^{-2i+2}\bigr).
\]
This contradicts $A>B$ and $i\ge2$. It proves
\eqref{eq:prime-coordinate-no-canonical-translation}, and hence
$\dim V_i\le1$ whenever that hypothesis holds. Equality in the
ramification estimate has not been asserted for a coordinate of
$p$-divisible valuation.

Write $A=p^aA_0$, with $p\nmid A_0$, and set
\[
 \delta_i=a+d_i-2i+2.
\]
The coordinate in \eqref{eq:proper-point-values} has positive
integer valuation $A_0p^{\delta_i}$ in $F_i$, so $\delta_i\ge0$.
Moreover
\[
 \delta_{i+1}-\delta_i=(d_{i+1}-d_i)-2\in\{-2,-1,0\}.
\]
If $2\le i<n$ and $\delta_i=0$, then integrality at the next
level forces $\delta_{i+1}=0$ and $d_{i+1}-d_i=2$.
Therefore $V_{i+1}=G[p]$ contains the canonical line, while
the coordinate of $Q_{i+1}$ has valuation prime to $p$.
This contradicts
\eqref{eq:prime-coordinate-no-canonical-translation}.
We have proved $\delta_i\ge1$ for every proper level
$2\le i<n$, which gives the asserted stronger lower bound.

At $i=n$, integrality gives $d_n\ge\max\{0,2n-2-a\}$.
If $d_n=2n-2-a\ge0$, then $\delta_n=0$, so
\eqref{eq:prime-coordinate-no-canonical-translation} gives the
claimed degree and translation restriction. Every point in the
statement extends to a compatible tower through level $n$.
The proof therefore covers all the stated point fields.
\end{proof}

Rational first-level torsion imposes a further constraint on these
point-field towers: their actual translation spaces can only increase.
This observation strengthens the bound at the top available level.

\begin{lemma}\label{lem:point-field-translation-persistence}
Let $p$ be odd, let $G$ be a height-two $p$-divisible group over a
field $L$ of characteristic zero, and suppose that $G[p]$ is
pointwise $L$-rational. Let $Q_1,\ldots,Q_n$ be compatible points
of exact orders $p,\ldots,p^n$, with $pQ_j=Q_{j-1}$, and put
$E_j=L(Q_j)$. Identify the relative Galois group
$\operatorname{Gal}(E_j/E_{j-1})$ with its translation subspace
$V_j\subset G[p]$. Then
\[
 V_2\subseteq V_3\subseteq\cdots\subseteq V_n.
\]
Only the relative extensions $E_j/E_{j-1}$ are asserted to be
Galois; the fields $E_j/L$ need not be Galois.
\end{lemma}

\begin{proof}
Every conjugate of $Q_j$ over $E_{j-1}$ differs from it by an
$L$-rational $p$-point. Thus $E_j/E_{j-1}$ is Galois, and the
translation map identifies its group with an $\mathbf F_p$-subspace
of $G[p]$.

Let $T=T_p(G)$ and choose a primitive $u\in T$ whose reduction
modulo $p^n$ represents $Q_n$. Its reduction modulo $p^j$ represents
$Q_j$. Write $\rho$ for the Galois action on $T$. Given
$v\in V_j$, choose a lift $\gamma$ of its relative automorphism
to the absolute Galois group of $E_{j-1}$, and put
$M=\rho(\gamma)-1$. Rationality of $G[p]$ and the fact that
$\gamma$ fixes $Q_{j-1}$ give
\[
 M(T)\subseteq pT,\qquad Mu\in p^{j-1}T,
 \qquad v=\frac{Mu}{p^{j-1}}\pmod{pT}.
\]
The binomial expansion of the single endomorphism $M$ gives
\[
 (\rho(\gamma)^p-1)u
 =pMu+\sum_{k=2}^{p-1}\binom pk M^ku+M^pu.
\]
For $2\le k\le p-1$, the corresponding summand belongs to
$p^{j+k-1}T\subseteq p^{j+1}T$. The last term belongs to
$p^{j+p-2}T\subseteq p^{j+1}T$, where oddness of $p$ is used.
Consequently
\[
 (\rho(\gamma^p)-1)u\equiv pMu\pmod{p^{j+1}T}.
\]
In particular $\gamma^p$ fixes $Q_j$, and its translation on
$Q_{j+1}$ is
\[
 \frac{(\rho(\gamma^p)-1)u}{p^j}
 \equiv\frac{Mu}{p^{j-1}}=v\pmod{pT}.
\]
Hence $v\in V_{j+1}$. This proves each asserted inclusion.
The lift $\gamma$ need not have finite order or fix all the
$(j-1)$st-level torsion; only the displayed hypotheses were used.
\end{proof}

\begin{corollary}\label{cor:proper-point-field-full-steps}
In the configuration of Proposition~\ref{prop:proper-canonical-point-fields},
suppose $n\ge a+2$, where $a=v_p(A)\ge1$.
For every compatible canonical point tower, each relative extension
$L(Q_j)/L(Q_{j-1})$ with $a+1\le j\le n$ has degree $p^2$.
In particular every canonical exact $p^n$-point $Q$ satisfies
\[
 [L(Q):L]=p^d,\qquad d\ge2n-1-a.
\]
\end{corollary}

\begin{proof}
Put $E_j=L(Q_j)$, $d_j=\log_p[E_j:L]$ and $i_j=\dim_{\mathbf F_p}V_j$.
The preceding proposition gives $E_1=L$ and
$d_j=\sum_{l=2}^j i_l$. The lemma makes $i_j$ nondecreasing.
The index $a+1$ is a proper level and is at least two, so the
proper-level bound gives
\[
 d_{a+1}\ge2(a+1)-1-a=a+1.
\]
If $i_{a+1}\le1$, all the $a$ terms in this sum would be at most
one, giving $d_{a+1}\le a$, a contradiction. Thus
$V_{a+1}=G[p]$, and every subsequent translation space is also
$G[p]$. All the stated relative extensions have degree $p^2$.
Finally the proper-level bound at $n-1$ gives
\[
 d_n=d_{n-1}+2\ge2n-3-a+2=2n-1-a.
\]
Every canonical exact $p^n$-point supplies such a tower by
successive multiplication, which proves the all-point assertion.
\end{proof}

\begin{proof}[Proof of Corollary~\ref{cor:boundary-point-field-degrees}]
Theorem~\ref{thm:first-boundary-normal-form} supplies the exchanged
central configuration, in both directions. Its first-level values are
\[
 A=\frac{p^3-1}{2}p^r,\qquad
 B=\frac{p-1}{2}p^{r+2},\qquad v_p(A)=r.
\]
Proposition~\ref{prop:proper-canonical-point-fields} gives the
stated proper-level bounds. Here $n=r+3$ and $a=r$, so
Corollary~\ref{cor:proper-point-field-full-steps} applies and yields
\[
 d_n\ge2(r+3)-1-r=r+5.
\]
It also shows that the relative steps at $j=r+1,r+2,r+3$ all
have degree $p^2$.

For $r=1$ and $n=4$, the earlier proper-level bounds and the
fibre-cardinality upper bounds already give $d_2=2$ and $d_3=4$;
the denominator bound alone leaves $d_4\in\{5,6\}$. If $d_4=5$,
the earlier top equality clause would give a cyclic final extension
of degree $p$ with noncanonical translation line. The new
translation-space inclusion makes the final relative step have
degree $p^2$, and therefore $d_4=6$. Equivalently the three total
degrees are exactly $p^2,p^4,p^6$. This excludes the previous
degree-$p^5$ possibility without using an exact break formula
at a coordinate of $p$-divisible valuation.
\end{proof}

These are necessary point-field restrictions. The remaining
degree-$p^6$ alternative when $r=1$ is neither constructed nor
excluded, and the boundary character index is not evaluated here.
